\documentclass[11pt]{article}

\usepackage[margin=1in]{geometry}
\usepackage{amsmath,amssymb}
\usepackage{booktabs}
\usepackage{enumitem}
\usepackage{hyperref}
\hypersetup{hidelinks}

\title{Scaling Book Exercises: Section 12, Quiz 2}
\author{Lucas Yan}
\date{}

\begin{document}
\maketitle

\noindent
\textbf{Quiz:} GPU Nodes\\
\textbf{Source scan:} \texttt{../handwritten/section\_12\_handwritten.pdf},
pages 3--4\\
\textbf{Reference:} \url{https://jax-ml.github.io/scaling-book/gpus/}

\section*{Question 1: Total bandwidth of an H100 node}

Four NVSwitches with 64 ports each and \(25\) GB/s per port provide
\[
4\cdot64\cdot25=\boxed{6.4\ \text{TB/s}}
\]
of switch capacity.  The eight GPUs can each egress only \(450\) GB/s:
\[
8\cdot450=\boxed{3.6\ \text{TB/s}}.
\]
The smaller value is the bottleneck, so the usable aggregate node bandwidth
is
\[
\boxed{3.6\ \text{TB/s}.}
\]

\section*{Question 2: Bisection bandwidth}

Any even bisection has four GPUs on each side.  Counting traffic in both
directions, all eight GPUs can egress across the cut:
\[
\boxed{W_{\mathrm{bisect}}=8(450\ \text{GB/s})=3.6\ \text{TB/s}.}
\]

\section*{Question 3: Eight-GPU AllGather}

Let an array have total size \(B\) bytes and let \(N=8\).  Each GPU begins
with \(B/N\) bytes and sends one such chunk in each of \(N-1\) ring steps.
With per-GPU egress bandwidth \(W=450\) GB/s,
\[
T_{\mathrm{AG}}
=\frac{(N-1)B}{NW}
=\frac{7B}{8(450\times10^9)}
\approx\frac{B}{450\times10^9}.
\]

For \(\mathrm{bf16}[D_X,F]\) with \(D=4096\) and \(F=65{,}536\),
\[
B=2DF
=2(4096)(65{,}536)
=536{,}870{,}912\text{ bytes}.
\]
Therefore
\[
\boxed{
T_{\mathrm{AG}}
=\frac{7(536{,}870{,}912)}
       {8(450\times10^9)}
\approx1.04\text{ ms}.}
\]
The large-\(N\) approximation gives \(1.19\) ms; latency and imperfect
bandwidth utilization can raise the observed time to roughly \(1.5\) ms.

\end{document}
